\documentclass[aps,preprint,showpacs,superscriptaddress]{revtex4-2}  %for double-spaced preprint

\usepackage{graphicx}  % needed for figures
\usepackage{dcolumn}   % needed for some tables
\usepackage{bm}        % for math
\usepackage{amssymb}   % for math
\usepackage[caption=false]{subfig}
\usepackage{xcolor}
%\usepackage{ulem}       % for \sout

\bibliographystyle{apsrev4-2}

\newcommand{\minerva}{MINERvA}
\newcommand{\qelike}{quasielastic-like}
\newcommand{\tune}{\minerva~tune v4.4.1}
\newcommand{\ccnopi}{CC0$\pi$}
\newcommand{\enuqe}{\ensuremath{E_{\nu,QE}}}
\newcommand{\pt}{\ensuremath{p_{t}}}
\newcommand{\pz}{\ensuremath{p_{||}}}
\newcommand{\eavail}{\ensuremath{E_{\rm\scriptstyle available}}}
\newcommand{\sumtp}{\ensuremath{\Sigma T_{p}}}
\newcommand{\qqqe}{\ensuremath{Q^{2}_{\rm\scriptstyle (QE)}}}
\newcommand{\qzeroqe}{\ensuremath{q_{0}^{\rm\scriptstyle (QE)}}}

\makeatletter




\begin{document}


\begin{center}
\textbf{\large Supplemental Materials}
\end{center}



\section{Presentation of all slices of the Triple Differential Cross Sections}
Figures \ref{fig:sup_pzptsumtp_bin1}-\ref{fig:sup_pzptsumtp_bin6} show the ratio of the triple differential cross section to the reference model in the \pt-\pz-\sumtp, grouped in bins of $\pz$.
Figures \ref{fig:sup_q0emusumtp_bin1}-\ref{fig:sup_q0emusumtp_bin6} present the triple differential cross section in $\qzeroqe$-$E_\mu$-\sumtp, grouped in bins of $E_\mu$. Figures \ref{fig:sup_tripdiff_bin0}-\ref{fig:sup_tripdiff_bin4} split the main text's Figure 1 into individual \pz~slices to help disentangle regions with overlapping graphics.

\begin{figure}[p]
    \centering
    \includegraphics[width = 0.5\linewidth]{nu-3d-xsec-comps-pzptrec_no2p2Tune-0_resfsi-0_qefis-1_resisi-0_2p2htunes-0_ratio-1-pz_bin_1.eps}
    %\includegraphics[width=0.5\linewidth]{nu-3d-xsec-comps-pz_bin_log_all3D_ratio.eps}
    \caption{Ratio of the differential cross-section in panels of $\pt$ for the peak bin $1.5<\pz<3.5$~GeV/c.  The predicted cross section in the reference model is broken down into different contributions.  Note that in the highest $\pt$ panels, above $0.85$~GeV/c, the ratio is shown scaled by $0.5$ and $0.2$ relative to the other panels.}
    \label{fig:sup_pzptsumtp_bin1}
\end{figure}
\begin{figure}[tp]
    \centering
    \includegraphics[width = 0.5\linewidth]{nu-3d-xsec-comps-pzptrec_no2p2Tune-0_resfsi-0_qefis-1_resisi-0_2p2htunes-0_ratio-1-pz_bin_2.eps}
    %\includegraphics[width=0.5\linewidth]{nu-3d-xsec-comps-pz_bin_log_all3D_ratio.eps}
    \caption{Ratio of the differential cross-section in panels of $\pt$ for the peak bin $3.5<\pz<4.5$~GeV/c.  The predicted cross section in the reference model is broken down into different contributions.  Note that in the highest $\pt$ panels, above $0.85$~GeV/c, the ratio is shown scaled by $0.5$ relative to the other panels.}
    \label{fig:sup_pzptsumtp_bin2}
\end{figure}
\begin{figure}[tp]
    \centering
    \includegraphics[width = 0.5\linewidth]{nu-3d-xsec-comps-pzptrec_no2p2Tune-0_resfsi-0_qefis-1_resisi-0_2p2htunes-0_ratio-1-pz_bin_3.eps}
    %\includegraphics[width=0.5\linewidth]{nu-3d-xsec-comps-pz_bin_log_all3D_ratio.eps}
    \caption{Ratio of the differential cross-section in panels of $\pt$ for the peak bin $4.5<\pz<7.0$~GeV/c.  The predicted cross section in the reference model is broken down into different contributions.  Note that in the highest $\pt$ panels, above $0.85$~GeV/c, the ratio is shown scaled by $0.5$ relative to the other panels.}
    \label{fig:sup_pzptsumtp_bin3}
\end{figure}
\begin{figure}[tp]
    \centering
    \includegraphics[width = 0.5\linewidth]{nu-3d-xsec-comps-pzptrec_no2p2Tune-0_resfsi-0_qefis-1_resisi-0_2p2htunes-0_ratio-1-pz_bin_4.eps}
    %\includegraphics[width=0.5\linewidth]{nu-3d-xsec-comps-pz_bin_log_all3D_ratio.eps}
    \caption{Ratio of the differential cross-section in panels of $\pt$ for the peak bin $7.0<\pz<8.0$~GeV/c.  The predicted cross section in the reference model is broken down into different contributions.  Note that in the highest $\pt$ panels, above $1.0$~GeV/c, the ratio is shown scaled by $0.5$ relative to the other panels.}
    \label{fig:sup_pzptsumtp_bin4}
\end{figure}
\begin{figure}[tp]
    \centering
    \includegraphics[width = 0.5\linewidth]{nu-3d-xsec-comps-pzptrec_no2p2Tune-0_resfsi-0_qefis-1_resisi-0_2p2htunes-0_ratio-1-pz_bin_5.eps}
    %\includegraphics[width=0.5\linewidth]{nu-3d-xsec-comps-pz_bin_log_all3D_ratio.eps}
    \caption{Ratio of the differential cross-section in panels of $\pt$ for the peak bin $8.0<\pz<10.0$~GeV/c.  The predicted cross section in the reference model is broken down into different contributions.  Note that in the highest $\pt$ panels, above $0.85$~GeV/c, the ratio is shown scaled by $0.5$ relative to the other panels.}
    \label{fig:sup_pzptsumtp_bin5}
\end{figure}
\begin{figure}[tp]
    \centering
    \includegraphics[width = 0.5\linewidth]{nu-3d-xsec-comps-pzptrec_no2p2Tune-0_resfsi-0_qefis-1_resisi-0_2p2htunes-0_ratio-1-pz_bin_6.eps}
    %\includegraphics[width=0.5\linewidth]{nu-3d-xsec-comps-pz_bin_log_all3D_ratio.eps}
    \caption{Ratio of the differential cross-section in panels of $\pt$ for the peak bin $10.0<\pz<20.0$~GeV/c.  The predicted cross section in the reference model is broken down into different contributions.  Note that in the highest $\pt$ panels, above $0.70$~GeV/c, the ratio is shown scaled by $0.5$ relative to the other panels.}
\label{fig:sup_pzptsumtp_bin6}
\end{figure}



\begin{figure}[tp]
    \centering
    \includegraphics[width= 0.5\linewidth]{nu-3d-xsec-comps-enuproxyE_no2p2Tune-0_resfsi-0_qefis-1_resisi-0_2p2htunes-0_ratio-0-pz-multiplier_bin_1.eps}
    \caption{Flux-averaged triple differential cross section, $d^3\sigma/dE_\mu d\qzeroqe d\sumtp$.  Note that scale factors are applied to each of the $\qzeroqe$ panels.}
    \label{fig:sup_q0emusumtp_bin1}
\end{figure}

\begin{figure}[tp]
    \centering
    \includegraphics[width= 0.5\linewidth]{nu-3d-xsec-comps-enuproxyE_no2p2Tune-0_resfsi-0_qefis-1_resisi-0_2p2htunes-0_ratio-0-pz-multiplier_bin_2.eps}
    \caption{Flux-averaged triple differential cross section, $d^3\sigma/dE_\mu d\qzeroqe d\sumtp$.  Note that scale factors are applied to each of the $\qzeroqe$ panels.}
    \label{fig:sup_q0emusumtp_bin2}
\end{figure}

\begin{figure}[tp]
    \centering
    \includegraphics[width= 0.5\linewidth]{nu-3d-xsec-comps-enuproxyE_no2p2Tune-0_resfsi-0_qefis-1_resisi-0_2p2htunes-0_ratio-0-pz-multiplier_bin_3.eps}
    \caption{Flux-averaged triple differential cross section, $d^3\sigma/dE_\mu d\qzeroqe d\sumtp$.  Note that scale factors are applied to each of the $\qzeroqe$ panels.}
    \label{fig:sup_q0emusumtp_bin3}
\end{figure}
\begin{figure}[tp]
    \centering
    \includegraphics[width= 0.5\linewidth]{nu-3d-xsec-comps-enuproxyE_no2p2Tune-0_resfsi-0_qefis-1_resisi-0_2p2htunes-0_ratio-0-pz-multiplier_bin_4.eps}
    \caption{Flux-averaged triple differential cross section, $d^3\sigma/dE_\mu d\qzeroqe d\sumtp$.  Note that scale factors are applied to each of the $\qzeroqe$ panels.}
    \label{fig:sup_q0emusumtp_bin4}
\end{figure}
\begin{figure}[tp]
    \centering
    \includegraphics[width= 0.5\linewidth]{nu-3d-xsec-comps-enuproxyE_no2p2Tune-0_resfsi-0_qefis-1_resisi-0_2p2htunes-0_ratio-0-pz-multiplier_bin_5.eps}
    \caption{Flux-averaged triple differential cross section, $d^3\sigma/dE_\mu d\qzeroqe d\sumtp$.  Note that scale factors are applied to each of the $\qzeroqe$ panels.}
    \label{fig:sup_q0emusumtp_bin5}
\end{figure}
\begin{figure}[tp]
    \centering
    \includegraphics[width= 0.5\linewidth]{nu-3d-xsec-comps-enuproxyE_no2p2Tune-0_resfsi-0_qefis-1_resisi-0_2p2htunes-0_ratio-0-pz-multiplier_bin_6.eps}
    \caption{Flux-averaged triple differential cross section, $d^3\sigma/dE_\mu d\qzeroqe d\sumtp$.  Note that scale factors are applied to each of the $\qzeroqe$ panels.}
    \label{fig:sup_q0emusumtp_bin6}
\end{figure}



%Split the 3D

\begin{figure*}[tp]
    \centering
    \includegraphics[width= 0.9\linewidth]{nu-3d-xsec-comps-pzptrec-pz_bin_log_all3D_bin-0.eps}
    \caption{Triple differential cross section with only a single \pz~slice (1.5-3.5 GeV/c) shown The complete set of \pz~slices is shown in the main text, Figure 1.}
    \label{fig:sup_tripdiff_bin0}
\end{figure*}

\begin{figure*}[tp]
    \centering
    \includegraphics[width= 0.9\linewidth]{nu-3d-xsec-comps-pzptrec-pz_bin_log_all3D_bin-1.eps}
    \caption{Triple differential cross section with only a single \pz~slice (3.5-4.5 GeV/c) shown The complete set of \pz~slices is shown in the main text, Figure 1.}
    \label{fig:sup_tripdiff_bin1}
\end{figure*}

\begin{figure*}[tp]
    \centering
    \includegraphics[width= 0.9\linewidth]{nu-3d-xsec-comps-pzptrec-pz_bin_log_all3D_bin-2.eps}
    \caption{Triple differential cross section with only a single \pz~slice (4.5-7.0 GeV/c) shown The complete set of \pz~slices is shown in the main text, Figure 1.}
    \label{fig:sup_tripdiff_bin2}
\end{figure*}

\begin{figure*}[tp]
    \centering
    \includegraphics[width= 0.9\linewidth]{nu-3d-xsec-comps-pzptrec-pz_bin_log_all3D_bin-3.eps}
    \caption{Triple differential cross section with only a single \pz~slice (7.0-8.0 GeV/c) shown The complete set of \pz~slices is shown in the main text, Figure 1.}
    \label{fig:sup_tripdiff_bin3}
\end{figure*}

\begin{figure*}[tp]
    \centering
    \includegraphics[width= 0.9\linewidth]{nu-3d-xsec-comps-pzptrec-pz_bin_log_all3D_bin-4.eps}
    \caption{Triple differential cross section with only a single \pz~slice (8.0-10.0 GeV/c) shown The complete set of \pz~slices is shown in the main text, Figure 1.}
    \label{fig:sup_tripdiff_bin4}
\end{figure*}

\begin{figure*}[tp]
    \centering
    \includegraphics[width= 0.9\linewidth]{nu-3d-xsec-comps-pzptrec-pz_bin_log_all3D_bin-5.eps}
    \caption{Triple differential cross section with only a single \pz~slice (10.0-20.0 GeV/c) shown The complete set of \pz~slices is shown in the main text, Figure 1.}
    \label{fig:sup_tripdiff_bin5}
\end{figure*}




\end{document}